
\documentclass{article}
\usepackage{colortbl}
\usepackage{makecell}
\usepackage{multirow}
\usepackage{supertabular}

\begin{document}

\newcounter{utterance}

\centering \large Interaction Transcript for game `matchit\_ascii', experiment `similar\_grid\_1', episode 1 with Qwen3.6{-}35B{-}A3B{-}FP8{-}without{-}reasoning.
\vspace{24pt}

{ \footnotesize  \setcounter{utterance}{1}
\setlength{\tabcolsep}{0pt}
\begin{supertabular}{c@{$\;$}|p{.15\linewidth}@{}p{.15\linewidth}p{.15\linewidth}p{.15\linewidth}p{.15\linewidth}p{.15\linewidth}}
    \# & $\;$A & \multicolumn{4}{c}{Game Master} & $\;\:$B\\
    \hline

    \theutterance \stepcounter{utterance}  
    & & \multicolumn{4}{p{0.6\linewidth}}{
        \cellcolor[rgb]{0.9,0.9,0.9}{
            \makecell[{{p{\linewidth}}}]{
                \texttt{\tiny{[P1$\langle$GM]}}
                \texttt{U neemt deel aan een gezamenlijk raadspel.} \\
\texttt{Het doel is om erachter te komen of dit raster en een andere afbeelding die alleen ik kan zien, hetzelfde zijn.} \\
\texttt{Beschrijf eerst uw raster. Daarna geef ik mijn beschrijving en kunnen we elkaar vragen stellen over de rasters om erachter te komen of ze hetzelfde zijn. Begin nu uw korte rasterbeschrijving met "BESCHRIJVING:" gevolgd door de beschrijving. Voeg niets anders toe. Dit is uw raster:} \\
\texttt{X X X X X} \\
\texttt{X X X X ▢} \\
\texttt{X X X ▢ ▢} \\
\texttt{X X ▢ ▢ ▢} \\
\texttt{X ▢ ▢ ▢ ▢} \\
            }
        }
    }
    & \\ \\

    \theutterance \stepcounter{utterance}  
    & \multicolumn{4}{p{0.6\linewidth}}{
        \cellcolor[rgb]{0.9,0.9,0.9}{
            \makecell[{{p{\linewidth}}}]{
                \texttt{\tiny{[P1$\rangle$GM]}}
                \texttt{BESCHRIJVING: Een 5x5 raster. De bovenste rij bestaat volledig uit X's. In de tweede rij staat een ▢ op de vijfde positie. In de derde rij staan ▢ op de vierde en vijfde positie. In de vierde rij staan ▢ op de derde, vierde en vijfde positie. In de vijfde rij staan ▢ op de tweede, derde, vierde en vijfde positie. De X's vormen een driehoekige vorm in de linkerbovenhoek.} \\
            }
        }
    }
    & & \\ \\

    \theutterance \stepcounter{utterance}  
    & & & \multicolumn{2}{p{0.3\linewidth}}{
        \cellcolor[rgb]{0.9,0.9,0.9}{
            \makecell[{{p{\linewidth}}}]{
                \texttt{\tiny{[GM$|$GM]}}
                \texttt{continue} \\
            }
        }
    }
    & & \\ \\

    \theutterance \stepcounter{utterance}  
    & & \multicolumn{4}{p{0.6\linewidth}}{
        \cellcolor[rgb]{0.9,0.9,0.9}{
            \makecell[{{p{\linewidth}}}]{
                \texttt{\tiny{[GM$\rangle$P2]}}
                \texttt{U neemt deel aan een gezamenlijk raadspel.} \\
\texttt{Het doel is om erachter te komen of dit raster en een andere afbeelding die alleen ik kan zien, hetzelfde zijn.} \\
\texttt{Beschrijf eerst uw raster. Daarna geef ik mijn beschrijving en kunnen we elkaar vragen stellen over de rasters om erachter te komen of ze hetzelfde zijn. Begin nu uw korte rasterbeschrijving met "BESCHRIJVING:" gevolgd door de beschrijving. Voeg niets anders toe. Dit is uw raster:} \\
\texttt{X X X X X} \\
\texttt{▢ X X X X} \\
\texttt{▢ ▢ X X X} \\
\texttt{▢ ▢ ▢ X X} \\
\texttt{▢ ▢ ▢ ▢ X} \\
            }
        }
    }
    & \\ \\

    \theutterance \stepcounter{utterance}  
    & & & \multicolumn{4}{p{0.6\linewidth}}{
        \cellcolor[rgb]{0.9,0.9,0.9}{
            \makecell[{{p{\linewidth}}}]{
                \texttt{\tiny{[GM$\langle$P2]}}
                \texttt{BESCHRIJVING: Een 5x5 raster. De bovenste rij bestaat volledig uit X's. De tweede rij begint met een leeg vakje (▢) gevolgd door vier X's. De derde rij begint met twee lege vakjes gevolgd door drie X's. De vierde rij begint met drie lege vakjes gevolgd door twee X's. De onderste rij bestaat uit vier lege vakjes gevolgd door één X. De X's vormen een driehoekige vorm in de rechterbovenhoek.} \\
            }
        }
    }
     \\ \\

    \theutterance \stepcounter{utterance}  
    & & & \multicolumn{2}{p{0.3\linewidth}}{
        \cellcolor[rgb]{0.9,0.9,0.9}{
            \makecell[{{p{\linewidth}}}]{
                \texttt{\tiny{[GM$|$GM]}}
                \texttt{continue} \\
            }
        }
    }
    & & \\ \\

    \theutterance \stepcounter{utterance}  
    & & \multicolumn{4}{p{0.6\linewidth}}{
        \cellcolor[rgb]{0.9,0.9,0.9}{
            \makecell[{{p{\linewidth}}}]{
                \texttt{\tiny{[GM$\rangle$P2]}}
                \texttt{Dit is mijn BESCHRIJVINGBESCHRIJVING: Een 5x5 raster. De bovenste rij bestaat volledig uit X's. In de tweede rij staat een ▢ op de vijfde positie. In de derde rij staan ▢ op de vierde en vijfde positie. In de vierde rij staan ▢ op de derde, vierde en vijfde positie. In de vijfde rij staan ▢ op de tweede, derde, vierde en vijfde positie. De X's vormen een driehoekige vorm in de linkerbovenhoek.} \\
\texttt{Stel nu een vraag om nieuwe aspecten van mijn raster te ontdekken die anders kunnen zijn dan jouw raster. Begin met "VRAAG:" en voeg niets anders toe.} \\
            }
        }
    }
    & \\ \\

    \theutterance \stepcounter{utterance}  
    & & & \multicolumn{4}{p{0.6\linewidth}}{
        \cellcolor[rgb]{0.9,0.9,0.9}{
            \makecell[{{p{\linewidth}}}]{
                \texttt{\tiny{[GM$\langle$P2]}}
                \texttt{VRAAG: Zijn de lege vakjes (▢) in jouw raster allemaal identiek leeg, of bevatten ze een ander symbool of patroon?} \\
            }
        }
    }
     \\ \\

    \theutterance \stepcounter{utterance}  
    & & & \multicolumn{2}{p{0.3\linewidth}}{
        \cellcolor[rgb]{0.9,0.9,0.9}{
            \makecell[{{p{\linewidth}}}]{
                \texttt{\tiny{[GM$|$GM]}}
                \texttt{continue} \\
            }
        }
    }
    & & \\ \\

    \theutterance \stepcounter{utterance}  
    & & \multicolumn{4}{p{0.6\linewidth}}{
        \cellcolor[rgb]{0.9,0.9,0.9}{
            \makecell[{{p{\linewidth}}}]{
                \texttt{\tiny{[P1$\langle$GM]}}
                \texttt{Dit is mijn BESCHRIJVINGBESCHRIJVING: Een 5x5 raster. De bovenste rij bestaat volledig uit X's. De tweede rij begint met een leeg vakje (▢) gevolgd door vier X's. De derde rij begint met twee lege vakjes gevolgd door drie X's. De vierde rij begint met drie lege vakjes gevolgd door twee X's. De onderste rij bestaat uit vier lege vakjes gevolgd door één X. De X's vormen een driehoekige vorm in de rechterbovenhoek.} \\
\texttt{VRAAG: Zijn de lege vakjes (▢) in jouw raster allemaal identiek leeg, of bevatten ze een ander symbool of patroon?Begin je antwoord met "ANTWOORD:" en voeg niets anders toe.} \\
            }
        }
    }
    & \\ \\

    \theutterance \stepcounter{utterance}  
    & \multicolumn{4}{p{0.6\linewidth}}{
        \cellcolor[rgb]{0.9,0.9,0.9}{
            \makecell[{{p{\linewidth}}}]{
                \texttt{\tiny{[P1$\rangle$GM]}}
                \texttt{ANTWOORD: Ja, alle lege vakjes (▢) zijn identiek leeg en bevatten geen ander symbool of patroon.} \\
            }
        }
    }
    & & \\ \\

    \theutterance \stepcounter{utterance}  
    & & & \multicolumn{2}{p{0.3\linewidth}}{
        \cellcolor[rgb]{0.9,0.9,0.9}{
            \makecell[{{p{\linewidth}}}]{
                \texttt{\tiny{[GM$|$GM]}}
                \texttt{continue} \\
            }
        }
    }
    & & \\ \\

    \theutterance \stepcounter{utterance}  
    & & \multicolumn{4}{p{0.6\linewidth}}{
        \cellcolor[rgb]{0.9,0.9,0.9}{
            \makecell[{{p{\linewidth}}}]{
                \texttt{\tiny{[P1$\langle$GM]}}
                \texttt{Stel nu een vraag om nieuwe aspecten van mijn raster te ontdekken die anders kunnen zijn dan jouw raster. Begin met "VRAAG:" en voeg niets anders toe.} \\
            }
        }
    }
    & \\ \\

    \theutterance \stepcounter{utterance}  
    & \multicolumn{4}{p{0.6\linewidth}}{
        \cellcolor[rgb]{0.9,0.9,0.9}{
            \makecell[{{p{\linewidth}}}]{
                \texttt{\tiny{[P1$\rangle$GM]}}
                \texttt{VRAAG: Zijn de X's in jouw raster allemaal identiek, of hebben ze een andere vorm, grootte of kleur dan de X's in mijn raster?} \\
            }
        }
    }
    & & \\ \\

    \theutterance \stepcounter{utterance}  
    & & & \multicolumn{2}{p{0.3\linewidth}}{
        \cellcolor[rgb]{0.9,0.9,0.9}{
            \makecell[{{p{\linewidth}}}]{
                \texttt{\tiny{[GM$|$GM]}}
                \texttt{continue} \\
            }
        }
    }
    & & \\ \\

    \theutterance \stepcounter{utterance}  
    & & \multicolumn{4}{p{0.6\linewidth}}{
        \cellcolor[rgb]{0.9,0.9,0.9}{
            \makecell[{{p{\linewidth}}}]{
                \texttt{\tiny{[GM$\rangle$P2]}}
                \texttt{ANTWOORD: Ja, alle lege vakjes (▢) zijn identiek leeg en bevatten geen ander symbool of patroon.} \\
\texttt{VRAAG: Zijn de X's in jouw raster allemaal identiek, of hebben ze een andere vorm, grootte of kleur dan de X's in mijn raster?Begin je antwoord met "ANTWOORD:" en voeg niets anders toe.} \\
            }
        }
    }
    & \\ \\

    \theutterance \stepcounter{utterance}  
    & & & \multicolumn{4}{p{0.6\linewidth}}{
        \cellcolor[rgb]{0.9,0.9,0.9}{
            \makecell[{{p{\linewidth}}}]{
                \texttt{\tiny{[GM$\langle$P2]}}
                \texttt{ANTWOORD: Ja, alle X's in mijn raster zijn identiek.} \\
            }
        }
    }
     \\ \\

    \theutterance \stepcounter{utterance}  
    & & & \multicolumn{2}{p{0.3\linewidth}}{
        \cellcolor[rgb]{0.9,0.9,0.9}{
            \makecell[{{p{\linewidth}}}]{
                \texttt{\tiny{[GM$|$GM]}}
                \texttt{continue} \\
            }
        }
    }
    & & \\ \\

    \theutterance \stepcounter{utterance}  
    & & \multicolumn{4}{p{0.6\linewidth}}{
        \cellcolor[rgb]{0.9,0.9,0.9}{
            \makecell[{{p{\linewidth}}}]{
                \texttt{\tiny{[GM$\rangle$P2]}}
                \texttt{Stel nu een vraag om nieuwe aspecten van mijn raster te ontdekken die anders kunnen zijn dan jouw raster. Begin met "VRAAG:" en voeg niets anders toe.} \\
            }
        }
    }
    & \\ \\

    \theutterance \stepcounter{utterance}  
    & & & \multicolumn{4}{p{0.6\linewidth}}{
        \cellcolor[rgb]{0.9,0.9,0.9}{
            \makecell[{{p{\linewidth}}}]{
                \texttt{\tiny{[GM$\langle$P2]}}
                \texttt{VRAAG: Zijn de X's in jouw raster ook allemaal identiek en op dezelfde posities geplaatst als in mijn beschrijving, of zijn er afwijkingen in de vorm of positie van de X's?} \\
            }
        }
    }
     \\ \\

    \theutterance \stepcounter{utterance}  
    & & & \multicolumn{2}{p{0.3\linewidth}}{
        \cellcolor[rgb]{0.9,0.9,0.9}{
            \makecell[{{p{\linewidth}}}]{
                \texttt{\tiny{[GM$|$GM]}}
                \texttt{continue} \\
            }
        }
    }
    & & \\ \\

    \theutterance \stepcounter{utterance}  
    & & \multicolumn{4}{p{0.6\linewidth}}{
        \cellcolor[rgb]{0.9,0.9,0.9}{
            \makecell[{{p{\linewidth}}}]{
                \texttt{\tiny{[P1$\langle$GM]}}
                \texttt{ANTWOORD: Ja, alle X's in mijn raster zijn identiek.} \\
\texttt{VRAAG: Zijn de X's in jouw raster ook allemaal identiek en op dezelfde posities geplaatst als in mijn beschrijving, of zijn er afwijkingen in de vorm of positie van de X's?Begin je antwoord met "ANTWOORD:" en voeg niets anders toe.} \\
            }
        }
    }
    & \\ \\

    \theutterance \stepcounter{utterance}  
    & \multicolumn{4}{p{0.6\linewidth}}{
        \cellcolor[rgb]{0.9,0.9,0.9}{
            \makecell[{{p{\linewidth}}}]{
                \texttt{\tiny{[P1$\rangle$GM]}}
                \texttt{ANTWOORD: Nee, de X's zijn niet op dezelfde posities geplaatst. In mijn raster bevinden de X's zich in de linkerbovenhoek, terwijl in jouw beschrijving de X's in de rechterbovenhoek staan.} \\
            }
        }
    }
    & & \\ \\

    \theutterance \stepcounter{utterance}  
    & & & \multicolumn{2}{p{0.3\linewidth}}{
        \cellcolor[rgb]{0.9,0.9,0.9}{
            \makecell[{{p{\linewidth}}}]{
                \texttt{\tiny{[GM$|$GM]}}
                \texttt{continue} \\
            }
        }
    }
    & & \\ \\

    \theutterance \stepcounter{utterance}  
    & & \multicolumn{4}{p{0.6\linewidth}}{
        \cellcolor[rgb]{0.9,0.9,0.9}{
            \makecell[{{p{\linewidth}}}]{
                \texttt{\tiny{[P1$\langle$GM]}}
                \texttt{Stel nu een vraag om nieuwe aspecten van mijn raster te ontdekken die anders kunnen zijn dan jouw raster. Begin met "VRAAG:" en voeg niets anders toe.} \\
            }
        }
    }
    & \\ \\

    \theutterance \stepcounter{utterance}  
    & \multicolumn{4}{p{0.6\linewidth}}{
        \cellcolor[rgb]{0.9,0.9,0.9}{
            \makecell[{{p{\linewidth}}}]{
                \texttt{\tiny{[P1$\rangle$GM]}}
                \texttt{VRAAG: Zijn de X's in jouw raster allemaal identiek en op dezelfde posities geplaatst als in mijn beschrijving, of zijn er afwijkingen in de vorm of positie van de X's?} \\
            }
        }
    }
    & & \\ \\

    \theutterance \stepcounter{utterance}  
    & & & \multicolumn{2}{p{0.3\linewidth}}{
        \cellcolor[rgb]{0.9,0.9,0.9}{
            \makecell[{{p{\linewidth}}}]{
                \texttt{\tiny{[GM$|$GM]}}
                \texttt{continue} \\
            }
        }
    }
    & & \\ \\

    \theutterance \stepcounter{utterance}  
    & & \multicolumn{4}{p{0.6\linewidth}}{
        \cellcolor[rgb]{0.9,0.9,0.9}{
            \makecell[{{p{\linewidth}}}]{
                \texttt{\tiny{[GM$\rangle$P2]}}
                \texttt{ANTWOORD: Nee, de X's zijn niet op dezelfde posities geplaatst. In mijn raster bevinden de X's zich in de linkerbovenhoek, terwijl in jouw beschrijving de X's in de rechterbovenhoek staan.} \\
\texttt{VRAAG: Zijn de X's in jouw raster allemaal identiek en op dezelfde posities geplaatst als in mijn beschrijving, of zijn er afwijkingen in de vorm of positie van de X's?Begin je antwoord met "ANTWOORD:" en voeg niets anders toe.} \\
            }
        }
    }
    & \\ \\

    \theutterance \stepcounter{utterance}  
    & & & \multicolumn{4}{p{0.6\linewidth}}{
        \cellcolor[rgb]{0.9,0.9,0.9}{
            \makecell[{{p{\linewidth}}}]{
                \texttt{\tiny{[GM$\langle$P2]}}
                \texttt{ANTWOORD: Nee, de X's zijn niet op dezelfde posities geplaatst. In mijn raster bevinden de X's zich in de rechterbovenhoek, terwijl in jouw beschrijving de X's in de linkerbovenhoek staan.} \\
            }
        }
    }
     \\ \\

    \theutterance \stepcounter{utterance}  
    & & & \multicolumn{2}{p{0.3\linewidth}}{
        \cellcolor[rgb]{0.9,0.9,0.9}{
            \makecell[{{p{\linewidth}}}]{
                \texttt{\tiny{[GM$|$GM]}}
                \texttt{continue} \\
            }
        }
    }
    & & \\ \\

    \theutterance \stepcounter{utterance}  
    & & \multicolumn{4}{p{0.6\linewidth}}{
        \cellcolor[rgb]{0.9,0.9,0.9}{
            \makecell[{{p{\linewidth}}}]{
                \texttt{\tiny{[GM$\rangle$P2]}}
                \texttt{Stel nu een vraag om nieuwe aspecten van mijn raster te ontdekken die anders kunnen zijn dan jouw raster. Begin met "VRAAG:" en voeg niets anders toe.} \\
            }
        }
    }
    & \\ \\

    \theutterance \stepcounter{utterance}  
    & & & \multicolumn{4}{p{0.6\linewidth}}{
        \cellcolor[rgb]{0.9,0.9,0.9}{
            \makecell[{{p{\linewidth}}}]{
                \texttt{\tiny{[GM$\langle$P2]}}
                \texttt{VRAAG: Zijn de X's in jouw raster allemaal identiek qua vorm en grootte, of zijn er variaties in de X's zelf?} \\
            }
        }
    }
     \\ \\

    \theutterance \stepcounter{utterance}  
    & & & \multicolumn{2}{p{0.3\linewidth}}{
        \cellcolor[rgb]{0.9,0.9,0.9}{
            \makecell[{{p{\linewidth}}}]{
                \texttt{\tiny{[GM$|$GM]}}
                \texttt{continue} \\
            }
        }
    }
    & & \\ \\

    \theutterance \stepcounter{utterance}  
    & & \multicolumn{4}{p{0.6\linewidth}}{
        \cellcolor[rgb]{0.9,0.9,0.9}{
            \makecell[{{p{\linewidth}}}]{
                \texttt{\tiny{[P1$\langle$GM]}}
                \texttt{ANTWOORD: Nee, de X's zijn niet op dezelfde posities geplaatst. In mijn raster bevinden de X's zich in de rechterbovenhoek, terwijl in jouw beschrijving de X's in de linkerbovenhoek staan.} \\
\texttt{VRAAG: Zijn de X's in jouw raster allemaal identiek qua vorm en grootte, of zijn er variaties in de X's zelf?Begin je antwoord met "ANTWOORD:" en voeg niets anders toe.} \\
            }
        }
    }
    & \\ \\

    \theutterance \stepcounter{utterance}  
    & \multicolumn{4}{p{0.6\linewidth}}{
        \cellcolor[rgb]{0.9,0.9,0.9}{
            \makecell[{{p{\linewidth}}}]{
                \texttt{\tiny{[P1$\rangle$GM]}}
                \texttt{ANTWOORD: Ja, de X's in mijn raster zijn allemaal identiek qua vorm en grootte.} \\
            }
        }
    }
    & & \\ \\

    \theutterance \stepcounter{utterance}  
    & & & \multicolumn{2}{p{0.3\linewidth}}{
        \cellcolor[rgb]{0.9,0.9,0.9}{
            \makecell[{{p{\linewidth}}}]{
                \texttt{\tiny{[GM$|$GM]}}
                \texttt{continue} \\
            }
        }
    }
    & & \\ \\

    \theutterance \stepcounter{utterance}  
    & & \multicolumn{4}{p{0.6\linewidth}}{
        \cellcolor[rgb]{0.9,0.9,0.9}{
            \makecell[{{p{\linewidth}}}]{
                \texttt{\tiny{[P1$\langle$GM]}}
                \texttt{Stel nu een vraag om nieuwe aspecten van mijn raster te ontdekken die anders kunnen zijn dan jouw raster. Begin met "VRAAG:" en voeg niets anders toe.} \\
            }
        }
    }
    & \\ \\

    \theutterance \stepcounter{utterance}  
    & \multicolumn{4}{p{0.6\linewidth}}{
        \cellcolor[rgb]{0.9,0.9,0.9}{
            \makecell[{{p{\linewidth}}}]{
                \texttt{\tiny{[P1$\rangle$GM]}}
                \texttt{VRAAG: Zijn de lege vakjes (▢) in jouw raster allemaal identiek qua vorm en grootte, of zijn er variaties in de lege vakjes zelf?} \\
            }
        }
    }
    & & \\ \\

    \theutterance \stepcounter{utterance}  
    & & & \multicolumn{2}{p{0.3\linewidth}}{
        \cellcolor[rgb]{0.9,0.9,0.9}{
            \makecell[{{p{\linewidth}}}]{
                \texttt{\tiny{[GM$|$GM]}}
                \texttt{continue} \\
            }
        }
    }
    & & \\ \\

    \theutterance \stepcounter{utterance}  
    & & \multicolumn{4}{p{0.6\linewidth}}{
        \cellcolor[rgb]{0.9,0.9,0.9}{
            \makecell[{{p{\linewidth}}}]{
                \texttt{\tiny{[GM$\rangle$P2]}}
                \texttt{ANTWOORD: Ja, de X's in mijn raster zijn allemaal identiek qua vorm en grootte.} \\
\texttt{VRAAG: Zijn de lege vakjes (▢) in jouw raster allemaal identiek qua vorm en grootte, of zijn er variaties in de lege vakjes zelf?Begin je antwoord met "ANTWOORD:" en voeg niets anders toe.} \\
            }
        }
    }
    & \\ \\

    \theutterance \stepcounter{utterance}  
    & & & \multicolumn{4}{p{0.6\linewidth}}{
        \cellcolor[rgb]{0.9,0.9,0.9}{
            \makecell[{{p{\linewidth}}}]{
                \texttt{\tiny{[GM$\langle$P2]}}
                \texttt{ANTWOORD: Ja, de lege vakjes (▢) in mijn raster zijn allemaal identiek qua vorm en grootte.} \\
            }
        }
    }
     \\ \\

    \theutterance \stepcounter{utterance}  
    & & & \multicolumn{2}{p{0.3\linewidth}}{
        \cellcolor[rgb]{0.9,0.9,0.9}{
            \makecell[{{p{\linewidth}}}]{
                \texttt{\tiny{[GM$|$GM]}}
                \texttt{continue} \\
            }
        }
    }
    & & \\ \\

    \theutterance \stepcounter{utterance}  
    & & \multicolumn{4}{p{0.6\linewidth}}{
        \cellcolor[rgb]{0.9,0.9,0.9}{
            \makecell[{{p{\linewidth}}}]{
                \texttt{\tiny{[GM$\rangle$P2]}}
                \texttt{Neem nu een beslissing. Wat denk je: zijn jouw raster en het andere beschreven raster hetzelfde? Schrijf "BESLISSING: hetzelfde raster" als je denkt dat het hetzelfde raster is of "BESLISSING: verschillende rasters" als je denkt dat het verschillende rasters zijn. Voeg niets anders toe.} \\
            }
        }
    }
    & \\ \\

    \theutterance \stepcounter{utterance}  
    & & & \multicolumn{4}{p{0.6\linewidth}}{
        \cellcolor[rgb]{0.9,0.9,0.9}{
            \makecell[{{p{\linewidth}}}]{
                \texttt{\tiny{[GM$\langle$P2]}}
                \texttt{BESLISSING: verschillende rasters} \\
            }
        }
    }
     \\ \\

    \theutterance \stepcounter{utterance}  
    & & & \multicolumn{2}{p{0.3\linewidth}}{
        \cellcolor[rgb]{0.9,0.9,0.9}{
            \makecell[{{p{\linewidth}}}]{
                \texttt{\tiny{[GM$|$GM]}}
                \texttt{continue} \\
            }
        }
    }
    & & \\ \\

    \theutterance \stepcounter{utterance}  
    & & & \multicolumn{2}{p{0.3\linewidth}}{
        \cellcolor[rgb]{0.9,0.9,0.9}{
            \makecell[{{p{\linewidth}}}]{
                \texttt{\tiny{[GM$|$GM]}}
                \texttt{success} \\
            }
        }
    }
    & & \\ \\

    \theutterance \stepcounter{utterance}  
    & & \multicolumn{4}{p{0.6\linewidth}}{
        \cellcolor[rgb]{0.9,0.9,0.9}{
            \makecell[{{p{\linewidth}}}]{
                \texttt{\tiny{[P1$\langle$GM]}}
                \texttt{ANTWOORD: Ja, de lege vakjes (▢) in mijn raster zijn allemaal identiek qua vorm en grootte.} \\
\texttt{Neem nu een beslissing. Wat denk je: zijn jouw raster en het andere beschreven raster hetzelfde? Schrijf "BESLISSING: hetzelfde raster" als je denkt dat het hetzelfde raster is of "BESLISSING: verschillende rasters" als je denkt dat het verschillende rasters zijn. Voeg niets anders toe.} \\
            }
        }
    }
    & \\ \\

    \theutterance \stepcounter{utterance}  
    & \multicolumn{4}{p{0.6\linewidth}}{
        \cellcolor[rgb]{0.9,0.9,0.9}{
            \makecell[{{p{\linewidth}}}]{
                \texttt{\tiny{[P1$\rangle$GM]}}
                \texttt{BESLISSING: verschillende rasters} \\
            }
        }
    }
    & & \\ \\

    \theutterance \stepcounter{utterance}  
    & & & \multicolumn{2}{p{0.3\linewidth}}{
        \cellcolor[rgb]{0.9,0.9,0.9}{
            \makecell[{{p{\linewidth}}}]{
                \texttt{\tiny{[GM$|$GM]}}
                \texttt{continue} \\
            }
        }
    }
    & & \\ \\

    \theutterance \stepcounter{utterance}  
    & & & \multicolumn{2}{p{0.3\linewidth}}{
        \cellcolor[rgb]{0.9,0.9,0.9}{
            \makecell[{{p{\linewidth}}}]{
                \texttt{\tiny{[GM$|$GM]}}
                \texttt{success} \\
            }
        }
    }
    & & \\ \\

\end{supertabular}
}

\end{document}
